In addition to the 3rd party development libraries built using the open source
toolchains mentioned in \S\ref{sec:3rdparty}, \OHPC{} also provides {\em
  optional} compatible builds for use with the compilers and MPI stack included
in newer versions of the \IntelR{}~oneAPI HPC Toolkit. These packages provide a
similar hierarchical user environment experience as other compiler and MPI
families present in \OHPC{}.

To take advantage of the available builds, install the compatibility packages
that integrate oneAPI-generated compiler and MPI modulefiles within the standard
\OHPC{} user environment. On Ubuntu these compatibility packages configure the
Intel APT repository and install the required oneAPI packages through APT.

% begin_ohpc_run
% ohpc_comment_header Install Intel oneAPI tools \ref{sec:3rdparty_intel}
% ohpc_command if [[ ${enable_intel_packages} -eq 1 ]];then
% ohpc_indent 5
\begin{lstlisting}[language=bash,keywords={},upquote=true,keepspaces]
# Enable Intel oneAPI and install OpenHPC compatibility packages
[sms](*\#*) (*\install*) intel-compilers-devel-ohpc intel-mpi-devel-ohpc
\end{lstlisting}
% ohpc_indent 0
% ohpc_command fi
% end_ohpc_run

\begin{center}
\begin{tcolorbox}[]
As noted in \S\ref{sec:master_customization}, the default installation path for
\OHPC{} (\texttt{/opt/ohpc/pub}) is exported over NFS from the {\em master} to the
compute nodes, but the \IntelR{} oneAPI HPC Toolkit packages install to a top-level path of
\texttt{/opt/intel}. To make the \IntelR{} compilers available to the compute
nodes one must add an additional NFS export for \texttt{/opt/intel} that is mounted
on desired compute nodes.
\end{tcolorbox}
\end{center}

\noindent The Ubuntu package manifest in Appendix~\ref{appendix:manifest} lists
the Intel-family library builds available for local selection.
